% =============================================================================
% modules/_template.tex — Reusable module documentation template
%
% INSTRUCTIONS:
%   1. Copy this file to modules/[module_name].tex
%   2. Replace every TODO with real content
%   3. Add \include{modules/[module_name]} to main.tex
%   4. Run through the review checklist in LATEX_DOC_STANDARD.md
%
% DO NOT modify the section structure or ordering.
% DO NOT remove any section, even if content is minimal.
% =============================================================================

% Language badge — e.g., Java 17 · Spring Boot, or TypeScript · React

\chapter{TODO: Module Display Name}
\label{ch:TODO-module-id}

\modulelang{TODO: Language · Framework}

% Brief subtitle shown below the chapter rule (optional — remove if not needed)
\textcolor{CraftGray}{\textit{TODO: One-line role in the system (e.g., "Core REST API — players, balances, transactions")}}

\bigskip

% =============================================================================
\section{Overview}
% =============================================================================
%
% GUIDELINES:
%   - 2–4 sentences only.
%   - Answer: what is this module, where does it sit, who consumes it?
%   - Do NOT list features here. Features belong in Responsibilities.
%
TODO: Describe what the module is and where it fits in the Craftalism system.
Identify its primary consumers (other services, administrators, game players).
Do not list features or implementation details.

% =============================================================================
\section{Responsibilities}
% =============================================================================
%
% GUIDELINES:
%   - Bullet list of concrete, operational behaviors.
%   - Each bullet starts with a verb (Validates, Issues, Caches, Exposes…).
%   - One sentence per bullet.
%   - This is NOT a feature list — describe what the module does operationally.
%
\begin{itemize}
  \item TODO: Verb + concrete behavior (e.g., "Validates incoming JWTs against the published JWKS.").
  \item TODO: ...
  \item TODO: ...
  \item TODO: ...
\end{itemize}

% =============================================================================
\section{Architecture and Design}
% =============================================================================
%
% GUIDELINES:
%   - Describe the internal structure and interaction with other services.
%   - Must include: layer/component descriptions, request or data flow,
%     and any non-obvious design decisions.
%   - Use \subsection for distinct sub-topics (layers, security model, etc.).
%   - Include a diagram if the structure benefits from visual representation.
%
TODO: Introductory paragraph describing the high-level design approach
(e.g., layered architecture, event-driven, stateless resource server, etc.).

\subsection{TODO: Layer or Component Name}

TODO: Description of this layer/component and its responsibility.

\subsection{TODO: Request or Data Flow}

TODO: Step-by-step description of the primary request path through this module.
Use an \texttt{enumerate} list for sequential steps.

\begin{enumerate}
  \item TODO: Step 1.
  \item TODO: Step 2.
  \item TODO: Step 3.
\end{enumerate}

% Optional: architecture diagram
% Uncomment and adapt the tikzpicture below, or use \includegraphics.
%
% \begin{figure}[h]
% \centering
% \begin{tikzpicture}[node distance=1.4cm and 2.0cm]
%   \node[service]  (a) {Component A};
%   \node[service, right=of a] (b) {Component B};
%   \draw[flow] (a) -- node[flowlabel, above]{action} (b);
% \end{tikzpicture}
% \caption{TODO: Figure caption}
% \label{fig:TODO-diagram-id}
% \end{figure}

% Optional: code example
% \begin{bashlist}[caption={TODO: Caption}]
% TODO: command or code snippet
% \end{bashlist}

% =============================================================================
\section{Tech Stack}
% =============================================================================
%
% GUIDELINES:
%   - Three columns: Category | Technology | Notes.
%   - Notes can be blank but the column must remain.
%   - Do not abbreviate technology names (write "Spring Boot 3.5", not "SB 3.5").
%
\rowcolors{2}{CraftLight}{white}
\begin{tabularx}{\linewidth}{@{} L{3.0cm} L{5.2cm} X @{}}
  \toprule
  {\tableheadstyle Category} &
  {\tableheadstyle Technology} &
  {\tableheadstyle Notes} \\
  \midrule
  Language    & TODO: e.g., Java 17      & TODO: or leave blank \\
  Framework   & TODO: e.g., Spring Boot 3.5 & \\
  Database    & TODO: e.g., PostgreSQL 18 & \\
  Build Tool  & TODO: e.g., Gradle Wrapper & \\
  Testing     & TODO: e.g., JUnit 5, Mockito & \\
  Packaging   & TODO: e.g., Docker multi-stage image & \\
  \bottomrule
\end{tabularx}

% =============================================================================
\section{Configuration}
% =============================================================================
%
% GUIDELINES:
%   - Four columns: Variable | Default | Required | Description.
%   - Mark required variables with \textbf{Yes}.
%   - Use em dash (---) for variables with no default.
%   - If the module has no configuration, replace the table with:
%     "This module has no external configuration."
%
\rowcolors{2}{CraftLight}{white}
\begin{tabularx}{\linewidth}{@{} L{4.2cm} L{2.0cm} C{1.6cm} X @{}}
  \toprule
  {\tableheadstyle Variable} &
  {\tableheadstyle Default} &
  {\tableheadstyle Required} &
  {\tableheadstyle Description} \\
  \midrule
  \code{TODO\_VAR\_NAME}  & ---                & \textbf{Yes} & TODO: Description. \\
  \code{TODO\_VAR\_NAME}  & \code{TODO\_value} & No           & TODO: Description. \\
  \bottomrule
\end{tabularx}

% =============================================================================
\section{Known Limitations}
% =============================================================================
%
% GUIDELINES:
%   - Each item describes one specific gap, bug, or incomplete feature
%     in the CURRENT codebase. Write in present tense.
%   - Do NOT include planned improvements here — those go in Roadmap.
%   - If there are no known limitations, write "None at this time."
%
\begin{itemize}
  \item TODO: Present-tense statement of a current limitation
        (e.g., "Transaction creation does not update balances atomically.").
  \item TODO: ...
\end{itemize}

% =============================================================================
\section{Roadmap}
% =============================================================================
%
% GUIDELINES:
%   - Action-oriented planned improvements ("Add...", "Integrate...", "Replace...").
%   - Do NOT duplicate Known Limitations — the roadmap is the intended solution,
%     not a restatement of the problem.
%   - If there are no planned improvements, write "No roadmap items at this time."
%
\begin{itemize}
  \item TODO: Action-oriented improvement
        (e.g., "Integrate transaction creation with atomic balance transfer.").
  \item TODO: ...
\end{itemize}
